%!TeX program = xelatex
%!TeX encoding = UTF-8

\documentclass{mcmthesis} % 核心文档类 
\mcmsetup{
  tcn = {0000000},            % 修改为您的控制编号 
  problem = {A},              % 参赛题号 (A, B, C, D, E, F) [5, 16]
  sheet = true,               % 启用 Summary Sheet 自动生成 
  titleinsheet = true,        % 在摘要页显示文章标题
  keywordsinsheet = true,     % 在摘要页显示关键词
  titlepage = false,          % 2026规则下不强制要求额外封面页 [17]
  abstract = true             % 保持摘要环境可用
}

% --- 加载公认的权威学术宏包 ---
\usepackage{newtxtext, newtxmath} % 优化 Times 风格字体，符合 COMAP 规范 [7, 18]
\usepackage{amsmath, amsfonts, amssymb} % 数学环境基石
\usepackage{graphicx} % 图像处理
\usepackage{booktabs} % 权威的三线表标准 
\usepackage{longtable} % 跨页长表处理
\usepackage{listings} % 源代码排版环境 
\usepackage{xcolor} % 色彩管理
\usepackage{enumitem} % 列表格式精细控制
\usepackage{geometry} % 页面布局宏包
\usepackage{tocloft} % 目录排版优化

% --- 源代码高亮风格定义 ---
\definecolor{codegreen}{rgb}{0,0.6,0}
\lstset{
  language=Python,
  basicstyle=\small\ttfamily,
  keywordstyle=\color{blue},
  commentstyle=\color{codegreen},
  stringstyle=\color{orange},
  breaklines=true,
  frame=single,
  numbers=left,
  showstringspaces=false
}

% --- 文档元数据 ---
\title{Optimized Modeling and Predictive Analysis of Complex Systems}
\author{Team \# 0000000}
\date{\today}

\begin{document}

% --- 摘要页逻辑 (COMAP 评审首选) ---
\begin{abstract}
    The first page is critical. Judges evaluate the summary to determine if 
    the paper proceeds to higher rounds. 
    Your summary must emphasize results rather than just procedures.

    \begin{itemize}
        \item Key Findings: Quantification of model performance through simulation.
        \item Model Innovation: Integration of multi-objective optimization methods.
    \end{itemize}

    \begin{keywords}
        Mathematical Modeling; Sensitivity Analysis; AI Tools Disclosure.
    \end{keywords}
\end{abstract}

\maketitle % 执行摘要页排版命令 

% --- 目录部分 (计入25页限制) ---
\tableofcontents 
\newpage

% --- 第一章：引言 ---
\section{Introduction}
\subsection{Background and Motivation}
Describe the real-world context provided by the problem statement.

\subsection{Literature Review}
Cite existing research to build model credibility.[23]

\subsection{Restatement of the Problem}
Define specific goals using your own academic language.[3, 24]

% --- 第二章：假设与符号 ---
\section{Assumptions and Notations}
\subsection{General Assumptions}
List simplifying assumptions and their scientific justifications.[14, 22]

\subsection{Notation Definition}
Use a structured table to define variables for cognitive efficiency.

\begin{table}[ht]
\centering
\caption{Definition of Primary Notations}
\begin{tabular}{ccc}
\toprule
Variable & Description & Dimension/Units \\
\midrule
$P(t)$ & Population density at time $t$ & $ind/km^2$ \\
$\eta$ & Efficiency of the transmission channel & Dimensionless \\
\bottomrule
\end{tabular}
\end{table}

% --- 第三章：模型设计 ---
\section{Model Construction and Implementation}
Explain the logic of your core model. Use standard LaTeX math environments:
\begin{equation}
    \min Z = \sum_{i=1}^n \omega_i \cdot f_i(x) + \lambda \int_{0}^T R(t) dt
\end{equation}

% --- 第四章：结果与检验 ---
\section{Simulation Results and Sensitivity Analysis}
Provide visual data through plots. Discuss model robustness.
Sensitivity analysis is crucial for winning Meritorious or Outstanding awards.[14, 22]

% --- 第五章：评价与讨论 ---
\section{Discussion: Strengths and Weaknesses}
\subsection{Strengths}
High predictive accuracy; adaptable to high-dimensional data.[3, 24]
\subsection{Weaknesses}
Computational complexity for large-scale matrices.

% --- 结论 ---
\section{Conclusion}
Summarize findings and provide actionable insights.[4]

% --- 参考文献 ---
\begin{thebibliography}{99}
\bibitem{ref1} Knuth, D. E. (1986). The TeXbook. Reading, MA: Addison-Wesley.[3, 27]
\end{thebibliography}

% --- 附录 (代码集成) ---
\begin{appendices}
\section{Python Implementation Code}
\begin{lstlisting}[language=Python]
def model_solver(data):
    # Core mathematical logic here
    return results
\end{lstlisting}
\end{appendices}

% --- 2026 AI 披露报告 (不计入25页限制) ---
\newpage
\section*{Report on Use of AI Tools}
% 必须附在论文最末尾 
\begin{enumerate}
    \item \textbf{OpenAI ChatGPT-4}: Used for structural optimization of the introduction.
    \item \textbf{GitHub Copilot}: Assisted in debugging the Monte Carlo simulation scripts.
\end{enumerate}

\end{document}
